\tituloI{METODOLOGÍA}

\tituloII{Método, tipo y alcance de la investigación}

\tituloIII{Método de investigación}
\parrafoIII{Se adopta el método experimental. Se manipularán las herramientas de IA generativa (variable independiente) y los escenarios físicos (caída libre, colisiones y fricción) para observar su efecto sobre la precisión matemática y el rendimiento computacional (variables dependientes).}

\tituloIII{Tipo de investigación}
\parrafoIII{Investigación aplicada y tecnológica, orientada a establecer marcos de validación para el código generado automáticamente en el área de la computación física. El enfoque es experimental, dado que se controlan y manipulan las condiciones de generación de código para medir sus efectos en las variables dependientes.}

\tituloIII{Alcance de la investigación}
\parrafoIII{El alcance de la investigación es explicativo, ya que no solo se busca describir los errores y diferencias del código generado por IA, sino también explicar las relaciones causales entre el uso de estas herramientas y la degradación de la precisión matemática o el rendimiento computacional. Se establecen comparaciones sistemáticas frente a modelos teóricos y motores físicos de estándar industrial como Box2D.}

\tituloIII{Diseño de la investigación}
\parrafoIII{Se emplea un diseño experimental puro (comparativo). Se comparará el desempeño de los algoritmos de simulación física generados por IA frente a dos referencias: (1) el modelo matemático teórico (ecuaciones cinemáticas) y (2) un motor físico estándar (Box2D). Las mediciones se realizarán bajo condiciones de hardware idénticas y controladas, ejecutando un total de 30 repeticiones por cada escenario físico definido.}

\tituloII{Materiales y métodos}

\tituloIII{Materiales}

\tituloIII{Métodos}
